% ============================================================
%  Paper A — Methods
%  Full methodological detail is provided in the companion
%  Paper C (Estaji, in preparation). This section summarises
%  the approaches used for the six analyses reported here.
% ============================================================

\paragraph{Data sources and corpus construction.}
The corpus comprises \num{4.9} million plant science publications indexed in
OpenAlex~\cite{priem2022openalex} and PubMed.
OpenAlex records span 1900--2024; PubMed records span 1990--2024.
Records were retrieved using a controlled vocabulary of plant-science
subject tags (OpenAlex concept hierarchy, level 1--2) combined with
PubMed MeSH terms (Plant Physiology, Plant Biology, Botany, Agronomy,
and 47 related headings).
Duplicate records were resolved by DOI and then by title--author--year
triple matching, and publication years were capped at 2024 with pre-1900
records removed.
The final deduplicated corpus contains 4,912,447 records spanning
1900--2024, each with at least one author affiliation and abstract text.
Field-growth and organism-share analyses use the 1990--2024 modern window
($n \approx 3.98$ million), for which coverage from both sources is dense;
the historical analyses (temporal placebo, concept-marriage epochs, and the
orphan-crop attention gap) draw on the full 1900--2024 back-catalogue.
Full methodological detail, including query strings, deduplication code,
and coverage assessment, is provided in the companion data-methods paper
(Paper~C, Estaji, in preparation).

\paragraph{Growth modelling.}
Annual publication counts were fitted with both exponential
($y = a e^{bt}$) and logistic
($y = K / (1 + e^{-r(t - t_0)})$) growth models using nonlinear least-squares
optimisation (SciPy \texttt{curve\_fit}, Levenberg--Marquardt algorithm).
Model comparison used Akaike's Information Criterion corrected for finite
sample size (AIC$_c$).
Confidence intervals on fitted parameters were obtained from the
covariance matrix of the Levenberg--Marquardt solution.

\paragraph{Sleeping beauty identification.}
We computed the sleeping-beauty index $B$ as defined by Ke et al.~\cite{ke2015sleeping}:
\[
  B = \sum_{t=0}^{t_a} \left( \frac{c_{t_a}}{t_a} \cdot t - c_t \right),
\]
where $c_t$ is the citation count in year $t$ after publication and $t_a$ is
the awakening year, defined as the year in which the three-year forward
citation rate first exceeds ten times the mean pre-awakening citation rate.
Papers were required to have at least ten years of citation history and at
least 50 total citations by the end of the corpus window.
Organism assignment used title and abstract keyword matching against a
curated species lexicon of 847 plant taxa.
Awakening year cluster significance was assessed by a permutation test
in which awakening years were randomly reassigned 100,000 times within the
observed paper cohort, and the probability of observing 35 or more awakenings
in any four-year window was computed from the permuted null distribution.

\paragraph{Method diffusion modelling.}
Methodological categories were defined by Boolean keyword queries on title,
abstract, and author-keyword fields (full query strings in Supplementary
Table~\ref{tab:S3}).
Each method's annual publication share was fitted with a logistic
S-curve, and the time from introduction year to the year of crossing 5\%
share was extracted as the diffusion speed metric.
Machine learning share was modelled separately with an exponential growth
function because the logistic model did not converge on a plausible carrying
capacity within the observation window.
Reported growth rates are the fitted exponential rate constant.

\paragraph{Concept marriage detection.}
Author keywords and MeSH terms were extracted from all records and
normalised to a controlled vocabulary using exact-match lookup against
the AGROVOC thesaurus~\cite{caracciolo2013} supplemented by manual curation of
the 500 most frequent terms.
For each pair of concept terms, we computed annual Jaccard co-occurrence
similarity:
\[
  J_{ij}(y) = \frac{|P_i(y) \cap P_j(y)|}{|P_i(y) \cup P_j(y)|},
\]
where $P_i(y)$ is the set of papers containing concept $i$ in year $y$.
Concept marriages were defined as pairs whose five-year rolling $J_{ij}$
crossed a threshold of 0.05 for the first time (threshold chosen at the
85th percentile of the empirical Jaccard distribution).
Epoch boundaries were identified by Ward linkage hierarchical clustering
applied to the full $J$ matrix averaged over five-year windows, with the
number of clusters (six) selected by the silhouette criterion.

\paragraph{Semantic embedding (SPECTER2).}
SPECTER2~\cite{singh2022specter2} embeddings were computed for all 2.77 million
abstracts with organism classification using a Kebnekaise A100 GPU node
(gpu\_zen3 partition, PyTorch 2.7.1).
Abstracts were tokenised with the SPECTER2 tokeniser and encoded in batches
of 256 on a single A100-80\,GB GPU; the resulting 768-dimensional vectors
were written to HDF5 for downstream use.
Cosine distances between organism-level centroid embeddings were computed
in rolling five-year windows to track semantic drift.

\paragraph{Zero-shot organism classification.}
Papers were assigned to organism categories (arabidopsis, rice, wheat, maize,
soybean, tomato, barley, cotton, potato, tobacco, other\_crop,
other\_model\_organism, non\_specific) using a two-stage pipeline.
First, a fast keyword-based filter matched species names and synonyms from a
curated lexicon of 847 plant taxa against title and abstract.
Second, papers not resolved by keyword matching were passed through a
zero-shot classification head fine-tuned on top of SPECTER2 embeddings,
trained with 13 organism labels.
Organism confidence scores take one of two values. A score of 1.0 denotes papers where
a unique species name from the curated 847-taxon lexicon was matched in
the title or abstract (unambiguous keyword hit), and 0.7 denotes papers where
two or more species names matched (multi-match ambiguity resolved by
selecting the most frequently mentioned taxon).
Papers with confidence~0.7 constitute 4.6\% of the classified corpus.
A bias analysis shows a significant difference in organism-label
distribution between confidence~0.7 and confidence~1.0 papers
($\chi^2 = 537{,}346$, $p < 0.001$), though temporal coverage shows only
minimal divergence (KS\,=\,0.023, $p < 0.001$).
The confidence~0.7 papers represent genuine multi-species mentions and
are retained to avoid systematic exclusion of interdisciplinary research.
The 2.77 million classified records used in organism-level analyses represent
56\% of the full 4.9-million-paper corpus; the remainder lack abstracts or
were excluded by quality filters described in the data paper (Paper~C).

\paragraph{BERTopic modelling.}
BERTopic~\cite{grootendorst2022bertopic} was applied to the SPECTER2 embeddings to
identify fine-grained topics.
HDBSCAN ($\mathtt{min\_cluster\_size} = 50$) was used for initial clustering,
yielding 1,632 topics after outlier removal.
Of the 2,770,088 classified papers, 1,491,147 (53.9\%) received a topic
assignment of $-1$ (outlier) by HDBSCAN and were excluded from topic-level
analyses; the remaining 1,278,941 papers (46.1\%) carry a valid topic
assignment.
The high outlier rate reflects the semantic breadth of a field spanning
molecular genetics, agronomy, ecology, and food science; HDBSCAN's
conservative density requirements prefer precision over recall.
Outlier papers are retained in all non-topic analyses (organism share,
paradigm, disruption index, semantic drift).
Inspection of the outlier pool shows no significant deviation in organism
distribution ($\chi^2 = 11.4$, $df = 12$, $p = 0.51$) or decade composition
($\chi^2 = 3.2$, $df = 3$, $p = 0.36$) from the clustered papers, indicating
that the exclusion is approximately random with respect to the variables
studied.
Topics were aggregated into 30 macro-themes by hierarchical merging using
Ward linkage on the topic centroid distances, with the number of macro-themes
selected by the silhouette criterion.
Topic trend classification (emerging / stable / declining) was determined by
fitting a linear trend to the annual publication share of each macro-theme
over 2015--2024 and applying threshold criteria ($|\hat{\beta}| > 0.003$
per year, $p < 0.05$, Bonferroni-corrected).
Shannon entropy of the macro-theme distribution was computed annually as
$H = -\sum_k p_k \log_2 p_k$ where $p_k$ is the publication share of
macro-theme $k$ in year $t$.
Bonferroni correction is applied to: (1)~the $1{,}632$ pairwise topic
trend tests (threshold $\alpha / 1{,}632$), (2)~the Dunn pairwise
organism CD comparisons (45 pairs, threshold $\alpha / 45$), and
(3)~the organism--topic enrichment matrix ($13 \times 30 = 390$ cells,
threshold $\alpha / 390$).
Remaining tests (growth model comparison, sleeping-beauty permutation,
chi-square paradigm $\times$ organism, and regression-based attention gap)
are single or non-multiplied tests and are reported at $\alpha = 0.05$
without correction.

\paragraph{Paradigm classification.}
Each paper was classified as applied or fundamental using a keyword
taxonomy applied to title and abstract.
Applied papers include those mentioning yield, improvement, breeding,
agronomy, food security, crop production, or resistance/tolerance to
biotic/abiotic stress in a clearly translational framing.
Fundamental papers include those focused on mechanism elucidation,
protein structure, gene regulation, or developmental biology without
agronomic targets.
The independence of paradigm and organism category was assessed with a
Pearson chi-square test on the $13 \times 2$ contingency table.

\paragraph{CD disruption index by organism.}
The CD disruption index~\cite{funk2017disruption} was computed at decade resolution
for each organism group using forward and backward citation windows from
OpenAlex.
$\mathrm{CD} = (n_i - n_j) / (n_i + n_j + n_k)$, where $n_i$ is the
number of citing papers that cite the focal paper but not its references,
$n_j$ is the number that cite both, and $n_k$ is the number that cite
only references.
Organism-level differences in mean CD were assessed with the
Kruskal--Wallis test; pairwise comparisons used the Dunn test with
Bonferroni correction.

\paragraph{Self-citation and collaboration rates.}
Self-citation rate per organism was computed as the fraction of within-paper
citation pairs where at least one author of the citing paper is also an
author of the cited paper, using author-identifier matching from OpenAlex.
International collaboration rate was the fraction of papers with at least
two author affiliations in different countries, computed over 2020--2024.

\paragraph{Research attention gap.}
Economic importance scores were constructed from three FAO~\cite{fao2023} metrics:
gross production value (USD, five-year average 2019--2023), harvested area
(hectares), and number of countries in which the crop is a dietary staple
(defined as $\geq 5\%$ of national caloric supply).
All three metrics were log-transformed and standardised before combination
as an unweighted mean.
Publication counts were computed over 2020--2024 to match the economic
window.
The research attention gap is the studentised residual from an ordinary
least-squares regression of log publication count on log economic importance
score, fit across 120 species with complete data.
Species with $|z| > 2.0$ were classified as significantly over- or
under-studied.

\paragraph{Use of large language models.}
AI coding and writing assistants (Anthropic Claude and OpenAI models) were used
to support analysis-pipeline development, code review, figure and
supplementary-material organisation, and manuscript editing. The human authors
designed the study, verified all analyses, statistics, and figures, and take
full responsibility for the content and scientific claims. Large language models
do not satisfy authorship criteria, are not authors, and generated none of the
underlying data or results.
